\documentclass[10pt,a4paper]{amsart}
\usepackage[margin=19mm]{geometry}
\usepackage{amsmath,amssymb,mathtools}
\usepackage[scr=rsfs,cal=euler]{mathalfa}
\usepackage{xfrac}
\usepackage[unicode,hidelinks,bookmarks=false]{hyperref}
\newcommand{\ie}{i.e.\ }
\newcommand{\cf}{cf.\ }
\newcommand{\resp}{respectively\ }
\newcommand{\Subsection}{\subsection}
\providecommand{\nobreakdash}{\nobreak\hbox{-}}
\setlength{\emergencystretch}{2em}
\multlinegap0pt
\usepackage{amssymb}
\usepackage{mathtools}
\usepackage{xcolor}
\usepackage{enumerate}
\usepackage{tikz}
\newtheorem{theorem}{Theorem}
\newtheorem{corollary}{Corollary}
\newtheorem{proposition}{Proposition}
\title{A linear test approach to global controllability of third- and fifth-order nonlinear dispersive equations}
\author{Debanjit Mondal}
\date{Source: arXiv:2504.17580v2 (CC BY 4.0)}
\begin{document}
\maketitle

\noindent\emph{Source and licence.} A linear test approach to global controllability of third- and fifth-order nonlinear dispersive equations. Version arXiv:2504.17580v2 (CC BY 4.0), distributed by arXiv under the Creative Commons Attribution 4.0 International licence, http://creativecommons.org/licenses/by/4.0/, as recorded in the arXiv licence metadata for this exact version and shown on the abstract page https://arxiv.org/abs/2504.17580v2. Full licence notice: \texttt{licenses/arxiv-2504.17580-CC-BY-4.0.txt}.\par\smallskip

\noindent\textit{Editorial scope.} Counted unit: the article's proposition prp\_existence (source0.tex lines 201--203). The article's complete original proof of it is delivered with it (source0.tex lines 1017--1026, one \texttt{proof} environment of 10 lines, which the article places away from the statement). No mathematics is written, repaired or re-derived: the statement and the proof are the article's own lines, in the article's own notation, and no retained line is substituted or re-flowed.  Retained uncounted as the source-local context the proof needs: lines 83--88; lines 124--138; lines 141--146; lines 186--191; lines 205--210.  Nothing was omitted from any retained span, so the package carries no omission record.  No retained line contains a citation, so the wrapper carries no bibliography and no bibliography entry.  No retained line contains a figure environment, an image-inclusion command or drawing code, so no image is delivered and the package declares no asset dependency.  Editorial formatting only, in addition: an AMS \texttt{amsart} preamble that restates the article's own declarations and macros used by the retained lines, namely the environments \texttt{theorem}, \texttt{corollary}, \texttt{proposition} on separate counters, together with the standard packages those lines need.  The retained environments print in this document as Theorem 1, Corollary 1, Proposition 1 in order of appearance, whereas the article prints the same items with the numbering of its own full document.  The complete native source of the article as distributed in the arXiv e-print for this exact version is bundled unchanged as \texttt{sources/177037/source0.tex}.
\begin{align}\label{ctrl_prblm}  
	\begin{dcases}  
		\partial_t u + (-1)^{j + 1} \partial_x^{2j + 1} u + \frac{1}{2} \partial_x(u^2) = \eta(t, x), \quad (t, x) \in (0, T) \times \mathbb{T}, \\  
		u(0, x) = u^0 (x), \quad x \in \mathbb{T},  
	\end{dcases}  
\end{align}  

\begin{theorem}\label{main_thm}
	For any \( (j, s) \in \{1, 2\} \times \mathbb{N} \), the equation \eqref{ctrl_prblm} is approximately controllable in $H^{s+ 2j + 1}_0(\mathbb{T})$ at small time using a control that takes values in \( \mathcal{H} \). That is, for any $T > 0, $ any initial state \( u^0 \in H^{s+ 2j + 1}_0(\mathbb{T}) \) and target state \( u^1 \in H^{s+ 2j + 1}_0(\mathbb{T}) \), there exists a sufficiently small \( \tau > 0 \) and a control function \( \eta_{\tau, j} \in L^2((0, T\tau); \mathcal{H}) \) such that a solution \( u \) of \eqref{ctrl_prblm} satisfies  
	\begin{align}\label{lim_u_Ttau}  
		u_j(T \tau)\footnotemark \to u^1 \quad \text{in } H^s(\mathbb{T}) \text{ as } \tau \to 0^+.  
	\end{align}  
	Furthermore, the control \( \eta_{\tau, j} \) can be expressed as  
	\begin{align}\label{form_of_ctrl}  
		\eta_{\tau, j} = \mathcal{C}_{\tau}(u^0, u^1) + \kappa_{\tau, j},  
	\end{align}  
	where  
	\[
	\mathcal{C}_{\tau}: H^s_0(\mathbb{T}) \times H^s_0(\mathbb{T}) \to L^2((0, T\tau); \mathcal{H})  
	\]  
	is a bounded linear operator with a finite-dimensional range, and \( \kappa_{\tau, j} \in L^2((0, T\tau); \mathcal{H}) \). Importantly, both \( \mathcal{C}_{\tau} \) and \( \kappa_{\tau, j} \) are independent of the specific choices of \( u_0 \) and \( u_1 \). Limit \eqref{lim_u_Ttau} is uniform with respect to $u^0$ and $u^1$ in a bounded set of $H^{s+ 2j + 1}_0(\mathbb{T}).$
\end{theorem}

\begin{corollary}\label{Coro_to_cnclude_main_result}  
	For any \( (j, s) \in \{1, 2\} \times \mathbb{N} \), the equation \eqref{ctrl_prblm} exhibits global approximate controllability in \( H^s_0(\mathbb{T}) \) at any given time \( T > 0 \). That is, for any \( \varepsilon > 0 \) and any initial and target states \( u^0, u^1 \in H^s_0(\mathbb{T}) \), there exists a control function \( \eta_j \in L^2((0, T); \mathcal{H}) \) such that the corresponding solution \( u \) of \eqref{ctrl_prblm}, satisfies  
	\[
	\|u_j(T) - u^1 \|_{s} \leq \varepsilon.
	\]  
\end{corollary}

\begin{align}\label{diff_form_ctrl_prblm}  
	\begin{dcases}  
		\partial_t u + \mathcal{L}_j u + \mathcal{B}(u) = f, \quad (t, x) \in (0, T) \times \mathbb{T}, \\  
		u(0, x) = u^0 (x), \quad x \in \mathbb{T},  
	\end{dcases}  
\end{align}  

\begin{proposition}\label{prp_stability}  
	Let \( s \in \mathbb{N} \), \( T > 0 \), and \( f \in L^2((0, T); H^s_0(\mathbb{T})) \). For any initial data \( u^0, v^0 \in H^s(\mathbb{T}) \) satisfying \( [u^0] = [v^0] \), there exists a constant \( C(\|u^0\|_s, \|v^0\|_s) > 0 \) such that  
	\begin{align}\label{stability}  
		\sup_{t \in [0, T]} \|\mathcal{R}_{j, t}(u^0, f) - \mathcal{R}_{j, t}(v^0, f)\|_s \leq C \Big(\|u^0 - v^0\|_s \Big), \quad \forall j \in \{1, 2\}.  
	\end{align}  
\end{proposition}

\begin{proposition}\label{prp_existence}
	Let $(j, s) \in \{1, 2\} \times \mathbb{N}$, and suppose $T > 0$, \( u^0 \in H^s(\mathbb{T}) \), and \( f \in L^2((0, T); H^s_0(\mathbb{T})) \) are given. Then the system \eqref{diff_form_ctrl_prblm} admits a unique solution in \( C([0, T]; H^s(\mathbb{T})) \).  
\end{proposition}

\begin{proof}[Proof of Corollary~\ref{Coro_to_cnclude_main_result}]
	By the density of \( H^{s + 2j + 1}_0(\mathbb{T}) \) in \( H^s_0(\mathbb{T}) \) with semi-global stability \eqref{stability}, Theorem \ref{main_thm} ensures that for any \( u^0, u^1 \in H^s_0(\mathbb{T}) \), there exists a control \( \tilde{\eta}_{\tau, j} \in L^2((0, T\tau); \mathcal{H}) \) such that  
	\[
	\mathcal{R}_{j, T\tau}(u^0, \tilde{\eta}_{\tau, j}) \to u^1 \quad \text{in } H^s \text{ as } \tau \to 0^+.
	\]
	
	Now, by continuity of the flow in time and stability with respect to initial data (Propositions \ref{prp_existence}–\ref{prp_stability}), there exist \( \mathfrak{r} > 0 \), \( t' > 0 \) such that any trajectory starting in \( B(u^1, \mathfrak{r}) \) remains \( \varepsilon \)-close to \( u^1 \) for time \( t' \).
	
	Choosing \( \tau_1 \) small so that \( \mathcal{R}_{j, T\tau_1}(u^0, \tilde{\eta}_{\tau_1, j}) \in B(u^1, \mathfrak{r}) \), we patch the control with zero beyond \( T\tau_1 \) to extend the trajectory. If \( T\tau_1 + t' \geq T \), we are done. Otherwise, we iterate this argument finitely many times until reaching time \( T \). Thus, the system is approximately controllable at time \( T \) in \( H^s_0(\mathbb{T}) \).
\end{proof}

\end{document}
